\section{Reference algorithms and worked finite instance}
\label{app:reference}

This appendix describes the exact finite algorithms in the artifact.  They are
included for reproducibility and testing, not as a substitute for the Lean
formalization.

\subsection{Demand closure pseudocode}

\begin{lstlisting}[language=Python]
def close(seed, transfer, capacity_sum):
    d = seed
    history = [d]
    for _ in range(capacity_sum + 1):
        nxt = join(d, transfer(d))
        if nxt == d:
            return d, history
        assert rank(nxt) > rank(d)
        d = nxt
        history.append(d)
    raise Error("rank bound violated")
\end{lstlisting}
The implementation validates coordinate names and levels.  For small shapes it
exhaustively checks every admissible candidate to verify leastness.  It checks
monotonicity on every comparable pair when that set is below a configurable
limit; threshold-rule monotonicity is otherwise justified symbolically by the
form of the rules.

\subsection{Strategy checker pseudocode}

For a positive bounded formula:
\begin{lstlisting}[language=Python]
check(atom A, c, env) = (c.tag == atom and A(env))
check(P and Q, c, env) =
    c.tag == and and check(P,c.left,env) and check(Q,c.right,env)
check(P or Q, c, env) =
    c.side selects a branch whose subcertificate checks
check(exists i<N. P, c, env) =
    0 <= c.witness < N and check(P,c.sub,env[i := c.witness])
check(forall i<N. P, c, env) =
    c has N branches and each branch i checks under env[i := i]
\end{lstlisting}
The builder chooses the leftmost true disjunct and least true existential
witness.  These canonical choices model the object-level bounded-search plan.

\subsection{Backward viability pseudocode}

\begin{lstlisting}[language=Python]
viable[H] = terminal layer
for h = H-1 down to 0:
    viable[h] = [x in X[h] | exists y in viable[h+1], edge(h,x,y)]

if start not in viable[0]: fail
path = [start]
for h = 0 .. H-1:
    path.append(least y in viable[h+1] with edge(h,path[h],y))
\end{lstlisting}
The object formalization replaces host lists and predicates by bounded codes,
truth bits, finite folds, and checked least search.

\subsection{Generated demand instance}

The demonstration's public seed requests an approximate value but not a named
witness.  The closure reaches:
\begin{center}
\begin{tabular}{@{}lr@{}}
\toprule
Coordinate & Final level \\
\midrule
query:value & 1 \\
query:witness & 0 \\
cert:branch & 0 \\
cert:trace & 0 \\
precision:root & 5 \\
precision:integral & 6 \\
guard:separation & 0 \\
guard:denominator & 1 \\
residual:tail\_modulus & 1 \\
residual:height\_bound & 0 \\
\bottomrule
\end{tabular}
\end{center}
Thus the generated endpoint needs a denominator guard and tail modulus but not
branch separation, trace replay, or a Diophantine height bound.  This is the
kind of cross-category pruning that motivates the unified lattice.

\subsection{Test inventory}

The eleven unit tests cover:
\begin{enumerate}
  \item a cyclic demand closure;
  \item exhaustive leastness for the small test lattice;
  \item exhaustive monotonicity over comparable pairs;
  \item query-specific certificate liveness;
  \item residual and guard reporting;
  \item rejection of out-of-range rules;
  \item bounded least search;
  \item valid strategy construction;
  \item rejection of a tampered strategy and of a false existential;
  \item deterministic trace checking;
  \item viable and nonviable layered relations.
\end{enumerate}
The build script reruns these tests before compiling the PDF.
